\apendice{Plan de Proyecto Software}

\section{Introducción}
El plan de proyecto software es el documento que define cómo se va a gestionar y ejecutar el desarrollo de un producto o sistema informático. Establece el alcance, los objetivos, el cronograma, los recursos, el presupuesto, los riesgos y los criterios de calidad, así como la organización del equipo y los mecanismos de seguimiento y control. En esencia, convierte una idea de software en un itinerario gestionable, con responsabilidades claras y procedimientos para controlar el progreso y los cambios.

Su importancia radica en que reduce la incertidumbre y aumenta la probabilidad de éxito académico y técnico. Permite justificar decisiones, anticipar problemas, y demostrar competencias en gestión de proyectos, ingeniería de requisitos, calidad y comunicación técnica.

En este documento se va a incluir la planificación temporal de las fases y tareas del proyecto. Para cada una se debe establecer una fecha de inicio y se debe estimar la fecha final. Para realizar esto hay que tener en cuenta el peso de cada tarea, los requisitos previos y las dependencias.

También, se va a valorar la viabilidad del proyecto. En concreto, la viabilidad económica y la legal. En la viabilidad económica se estiman los costes y los beneficios del proyecto. Mientras que en la viabilidad legal analizan las leyes y licencias que afecten al desarrollo.  

\section{Planificación temporal}
Para la planificación se utiliza una metodología Scrum, aplicando iteraciones (\textit{sprints}). Al finalizar cada \textit{sprint} se realiza una reunión para la revisión del incremento y la planificación del siguiente. En la planificación se seleccionan y valoran las tareas (\textit{issues}) a realizar en el siguiente \textit{sprint}. Estas tareas se han valorado usando los \textit{story points} de ZenHub. 

\subsection{Sprint 0 (11/09/2025 - 18/09/2025)}
Este \textit{sprint} marco el comienzo del proyecto. Los objetivos fueron realizar una primera toma de contacto con los RAG y los LLM, investigando sobre que son, como funcionan y que ventajas aportan los RAG. También, se planteó empezar a escribir los conceptos teóricos relevantes y crear el repositorio \href{https://github.com/lbr1014/TFG_RAG.git}{GitHub} para el proyecto. Como el \textit{sprint} se diseño antes de la creación del repositorio no se creó el \textit{milestone} asociado como se ha hecho en los \textit{sprints} posteriores. 
Para llevar a cabo estos objetivos se plantearon las siguientes tareas:
\begin{itemize}
\tightlist
	\item Crear el repositorio en GitHub.
	\item Investigar sobre los RAG y LLM.
	\begin{itemize}
	\tightlist
		\item Leer el capitulo 1 del libro: A simple guide to retrieval augmented generation.
		\item Leer el capitulo 2 del libro: A simple guide to retrieval augmented generation.
	\end{itemize}
	\item Escribir conceptos teóricos sobre lo investigado.
\end{itemize}
Para realizar estas tareas se estimaron 6 horas de trabajo pero en realidad fueron 9. Al finalizar el \textit{sprint} se completaron todas las tareas.

\subsection{Sprint 1 (18/09/2025 - 25/09/2025)}
Los objetivos del \href{https://github.com/lbr1014/TFG_RAG/milestone/1}{Sprint 1} fueron investigar sobre los LLM y los RAG y escribir la información más relevante como concepto teórico en la memoria. Comenzar a escribir los objetivos del proyecto y las técnicas y herramientas definidas hasta el momento, como la metodología Scrum o el repositorio. 
Para llevar a cabo estos objetivos se plantearon las siguientes tareas:
\begin{itemize}
\tightlist
	\item Escribir objetivos del proyecto y técnicas y herramientas.
	\item Escribir conceptos teóricos.
	\begin{itemize}
	\tightlist
		\item Investigar sobre los RAG y los LLM.
		\begin{itemize}
		\tightlist
			\item Leer el capitulo 3 del libro: A simple guide to retrieval augmented generation.
			\ item Leer el capitulo 4 del libro: A simple guide to retrieval augmented generation.
		\end{itemize}
	\end{itemize}
\end{itemize}

Para realizar estas tareas se estimaron 7 horas de trabajo pero en realidad fueron 10. Al finalizar el \textit{sprint} se completaron todas las tareas.

El \textit{burndown} de este \textit{sprint} es:
\imagen{Burndown/Sprint1}{Sprint 1.}

\subsection{Sprint 2 (25/09/2025 - 02/10/2025)}
Los objetivos del \href{https://github.com/lbr1014/TFG_RAG/milestone/2}{Sprint 2} fueron seguir investigando sobre los LLM y RAG y explicar la información relevante en los conceptos teóricos.Investigar sobre el \textit{web scraping} y explicar en los anexos el plan de proyecto, en concreto la planificación temporal y los \textit{sprints}.
Para llevar a cabo estos objetivos se plantearon las siguientes tareas:
\begin{itemize}
\tightlist
	\item Añadir las correcciones a la memoria.
	\item Completar el apartado pipeline de generación del apartado conceptos teóricos.
	\item Investigar sobre web Scraping.
	\begin{itemize}
	\tightlist
		\item Investigar sobre HTML.
		\item Investigar sobre CSS.
	\end{itemize}
	\item Explicar el plan de proyectos en los anexos.
\end{itemize}

\section{Estudio de viabilidad}

\subsection{Viabilidad económica}

\subsection{Viabilidad legal}


